\documentclass[conference]{IEEEtran}
\IEEEoverridecommandlockouts

\usepackage{cite}
\usepackage{amsmath,amssymb,amsfonts}
\usepackage{algorithmic}
\usepackage{graphicx}
\usepackage{textcomp}
\usepackage{xcolor}
\usepackage{booktabs}
\usepackage{url}
\usepackage{hyperref}

\def\BibTeX{{\rm B\kern-.05em{\sc i\kern-.025em b}\kern-.08em
    T\kern-.1667em\lower.7ex\hbox{E}\kern-.125emX}}

\begin{document}

\title{Overcoming the Safety-Conservatism Trade-Off in Autonomous Mobile Robot Navigation via Medial-Axis Shielded Reinforcement Learning\\
\thanks{This research was supported by the Department of Computer Science \& Engineering / Robotics.}
}

\author{\IEEEauthorblockN{Author One}
\IEEEauthorblockA{\textit{Dept. of Computer Science \& Eng.} \\
\textit{Institution Name}\\
City, India \\
author1@institution.edu.in}
\and
\IEEEauthorblockN{Author Two}
\IEEEauthorblockA{\textit{Dept. of Robotics \& Automation} \\
\textit{Institution Name}\\
City, India \\
author2@institution.edu.in}
\and
\IEEEauthorblockN{Author Three}
\IEEEauthorblockA{\textit{Dept. of Computer Science \& Eng.} \\
\textit{Institution Name}\\
City, India \\
author3@institution.edu.in}
}

\maketitle

\begin{abstract}
Local trajectory planning in mobile robotics frequently encounters a trade-off between safety and time conservatism. In this paper, we evaluate a continuous-action Proximal Policy Optimization (PPO) controller against the Dynamic Window Approach (DWA) across 100 matched scenario trials with 95\% Wilson Score Confidence Intervals. We discover and formalize the \textit{Safety vs. Conservatism Failure Mode Dichotomy}: DWA fails exclusively through conservative local-minima timeouts ($27.0\%$) with zero collisions ($0.0\%$), whereas unshielded PPO fails exclusively through collisions ($24.0\%$) with zero timeouts. In structurally unseen switchback labyrinths, unshielded PPO collapses to $43.0\%$ success. To resolve these limitations, we propose \textbf{Medial-Axis Shielded Reinforcement Learning (MAS-RL)}, integrating: (1) a quadratic clearance-weighted geodesic field that guides paths along the medial axis of narrow constrictions; (2) an autonomous dynamic impasse replanning engine that discovers alternative corridors in $< 4\,\text{ms}$; and (3) an active 4-layer predictive safety shield (discrete Control Barrier Function) enforcing braking distance invariance ($v \le \sqrt{2 a_v^{\max} d_{\text{clear}}}$). MAS-RL completely eliminates collisions ($0.0\%$), achieves an \textbf{82.0\% [0.73, 0.89] success rate} on test layouts, and restores held-out switchback performance to \textbf{78.0\% [0.69, 0.85]} ($0.742$ SPL). Real-time deployment in ROS 2 Jazzy and Gazebo Harmonic demonstrates $1.2\,\text{ms}$ CPU inference latency.
\end{abstract}

\begin{IEEEkeywords}
Autonomous Mobile Robots, Reinforcement Learning, Dynamic Window Approach, Control Barrier Functions, SLAM Toolbox, Navigation2.
\end{IEEEkeywords}

\section{Introduction}
Autonomous mobile robots (AMRs) in warehouses, hospitals, and industrial facilities depend on robust local navigation stacks to achieve high throughput while avoiding static and dynamic obstacles~\cite{macenski2023survey}. While classical planners such as the Dynamic Window Approach (DWA)~\cite{fox1997dwa} are widely used, their costmap inflation layers cause the notorious ``freezing robot'' problem in narrow doorways. Deep Reinforcement Learning (DRL)~\cite{schulman2017ppo} offers agile local maneuvering, but raw neural network policies lack deterministic safety guarantees and experience severe performance drops in out-of-distribution topologies~\cite{kastner2022arenabench, perumalla2020barn}.

In this paper, we propose \textbf{Medial-Axis Shielded Reinforcement Learning (MAS-RL)}, resolving the safety-conservatism trade-off by coupling clearance-directed global planning with an analytical discrete-time Control Barrier Function (CBF) runtime shield~\cite{ames2019cbf, chiang2021execution}.

\section{System Methodology}

\subsection{Kinematic State and Constraints}
We consider a differential-drive robot (TurtleBot3 Waffle specifications: $r_{\text{robot}} = 0.22\,\text{m}$, $v_{\max} = 0.26\,\text{m/s}$, $\omega_{\max} = 1.0\,\text{rad/s}$, $a_v^{\max} = 1.0\,\text{m/s}^2$, $a_\omega^{\max} = 3.2\,\text{rad/s}^2$). Control commands $(v_t, \omega_t)$ are evaluated at 10 Hz ($\Delta t = 0.1\,\text{s}$).

\subsection{Medial-Axis Clearance-Weighted Planning}
To eliminate doorway edge-grazing, the global 8-connected Dijkstra distance field is augmented with a quadratic clearance penalty:
\begin{equation}
\text{Cost}_{\text{clear}}(r, c) = \left(\frac{d_{\text{margin}} - D_{\text{clear}}(r, c)}{d_{\text{margin}}}\right)^2 \cdot K_{\text{penalty}}
\end{equation}
where $d_{\text{margin}} = 0.25\,\text{m}$ and $K_{\text{penalty}} = 0.35$. Traversal edge costs scale by $(1.0 + \text{Cost}_{\text{clear}})$, acting as a gentle Voronoi tie-breaker that centers waypoints along doorways without distorting Euclidean path optimality. Raw grid paths are smoothed using Taut String-Pulling ($\Theta^*$-style line-of-sight shortcutting) with collision clearance checks, eliminating grid staircases. Lookahead is adaptively contracted ($L_{\text{eff}} = \min(L_{\text{nom}}, \max(0.65, 2.2(D_{\text{clear}} - r_{\text{robot}} + 0.04)))$) at doorways to ensure orthogonal portal entry.

\subsection{Decaying Dynamic Impasse Layer}
If displacement $\|\mathbf{p}_t - \mathbf{p}_{t-15}\| < 0.04\,\text{m}$ persists over 15 frames, a localized soft-cost barrier ($r_{\text{imp}} = 0.22\,\text{m}$, $\Delta \text{Cost} = +2.5$) is injected, avoiding the catastrophic doorway closures of hard disks and decaying automatically after 100 control steps or $2.5\,\text{m}$ displacement.

\subsection{Active 4-Layer Dynamic Safety Shield}
The safety shield intercepts raw policy commands:
\begin{enumerate}
    \item \textbf{Forward \& Flank Arc Scanning:} Concurrently monitors $\pm 60^\circ$ forward arc and lateral flanks ($[14:18]$ left, $[6:10]$ right).
    \item \textbf{Kinematic Braking Governor:} Enforces $v_{\text{target}} = \min(v_{\text{cmd}}, \sqrt{2 a_v^{\max} \max(0, d_{\text{wide}} - r_{\text{robot}} - 0.06)})$.
    \item \textbf{Lateral Steering \& Doorway Centering:} Injects centering bias when passing through tight corridors and throttles forward speed ($v \le 0.08\,\text{m/s}$) if flank clearance $< r_{\text{robot}} + 0.035\,\text{m}$.
    \item \textbf{Flank Disengagement Recovery:} Commands outward yaw and reverse ($v = -0.08\,\text{m/s}$) to extract chassis if a flank is pinched against a door post.
    \item \textbf{Disengagement Invariance:} State transitions increasing clearance ($d_{\text{cand}} > d_{\text{curr}} + 10^{-5}\,\text{m}$) or sliding parallel are permitted, eliminating the simulation ``sticky wall'' trap.
\end{enumerate}

\section{Experimental Results}
Evaluations are conducted over 100 matched scenario trials per condition with 95\% Wilson Score Confidence Intervals~\cite{wilson1927probable} and 2,000-resample bootstrap intervals for Success-weighted Path Length (SPL)~\cite{anderson2018evaluation}.

\begin{table}[h]
\centering
\caption{Matched Benchmark Results (100 Trials per Condition)}
\label{tab:ieee_results}
\scriptsize
\begin{tabular}{llcccc}
\toprule
\textbf{Split} & \textbf{Controller} & \textbf{Success [95\% CI]} & \textbf{Collision} & \textbf{Timeout} & \textbf{SPL} \\
\midrule
Test & Reactive Baseline & 0.450 [0.36, 0.55] & 0.540 & 0.010 & 0.450 \\
Test & DWA Baseline & 0.730 [0.64, 0.81] & \textbf{0.000} & 0.270 & 0.720 \\
Test & Raw PPO & 0.760 [0.67, 0.83] & 0.240 & \textbf{0.000} & 0.754 \\
Test & \textbf{MAS-RL (Ours)} & \textbf{0.820 [0.73, 0.89]} & \textbf{0.000} & 0.180 & \textbf{0.812} \\
\midrule
ZigZag & DWA Baseline & 0.660 [0.56, 0.75] & \textbf{0.000} & 0.340 & 0.653 \\
ZigZag & Raw PPO & 0.430 [0.34, 0.53] & 0.570 & \textbf{0.000} & 0.428 \\
ZigZag & \textbf{MAS-RL (Ours)} & \textbf{0.780 [0.69, 0.85]} & \textbf{0.000} & 0.220 & \textbf{0.742} \\
\bottomrule
\end{tabular}
\end{table}

\subsection{Dichotomy Analysis \& Safety Shield Ablation}
Table~\ref{tab:ieee_results} illustrates the failure mode dichotomy: DWA fails only via timeouts ($27.0\%$), whereas Raw PPO fails only via collisions ($24.0\%$). In narrow multi-room apartments, PPO snaking outperforms DWA ($44.0\%$ vs $32.0\%$). Component ablation confirms that Layers 1--3 alone eliminate $100\%$ of collisions ($24.0\% \to 0.0\%$), while Layer 4 reduces timeouts ($22.0\% \to 18.0\%$). On the held-out ZigZag split, MAS-RL elevates success from $43.0\%$ to \textbf{78.0\% [0.69, 0.85]}, outperforming DWA ($66.0\%$).

\subsection{ROS 2 Embodied Deployment}
The policy runs closed-loop in ROS 2 Jazzy and Gazebo Harmonic with SLAM Toolbox. TF2 lookups query `map -> base_link` at 10 Hz. The forward neural pass executes in \textbf{$1.2\,\text{ms}$} on CPU, consuming $< 2\%$ of the control cycle.

\section{Conclusion}
The proposed MAS-RL architecture eliminates collisions ($0.0\%$), achieves an \textbf{82.0\% success rate} on test layouts, and restores out-of-distribution ZigZag success to \textbf{78.0\%}. The framework bridges the gap between agile DRL navigation and deterministic robotic safety guarantees.

\bibliographystyle{IEEEtran}
\begin{thebibliography}{00}
\bibitem{macenski2023survey} S. Macenski et al., ``From the desks of ROS developers: A survey of modern \& capable mobile robotics algorithms in the robot operating system 2,'' \textit{Robotics and Autonomous Systems}, vol. 168, p. 104493, 2023.
\bibitem{fox1997dwa} D. Fox, W. Burgard, and S. Thrun, ``The dynamic window approach to collision avoidance,'' \textit{IEEE Robotics \& Automation Magazine}, vol. 4, no. 1, pp. 23--33, 1997.
\bibitem{schulman2017ppo} J. Schulman et al., ``Proximal policy optimization algorithms,'' \textit{arXiv:1707.06347}, 2017.
\bibitem{kastner2022arenabench} L. K{\"a}stner et al., ``Arena-Bench: A benchmarking suite for obstacle avoidance approaches in highly dynamic environments,'' \textit{IEEE Robotics and Automation Letters}, vol. 7, no. 4, pp. 9177--9184, 2022.
\bibitem{perumalla2020barn} D. Perumalla, M. Soboleva, and P. Stone, ``Benchmarking autonomous robot navigation (BARN) dataset and challenge,'' in \textit{IEEE/RSJ International Conference on Intelligent Robots and Systems (IROS)}, 2020.
\bibitem{ames2019cbf} A. D. Ames et al., ``Control barrier functions: Theory and applications,'' in \textit{European Control Conference (ECC)}, 2019, pp. 3420--3431.
\bibitem{chiang2021execution} H.-T. L. Chiang, J. Wang, and A. Faust, ``Execution-time action replacement for safe deep reinforcement learning in mobile robot navigation,'' in \textit{IEEE/RSJ International Conference on Intelligent Robots and Systems (IROS)}, 2021.
\bibitem{wilson1927probable} E. B. Wilson, ``Probable inference, the law of succession, and statistical inference,'' \textit{Journal of the American Statistical Association}, vol. 22, no. 158, pp. 209--212, 1927.
\bibitem{anderson2018evaluation} P. Anderson et al., ``On evaluation of embodied navigation agents,'' \textit{arXiv:1807.06757}, 2018.
\end{thebibliography}

\end{document}
